\documentclass[a4paper,twoside]{article}

\usepackage{graphicx}  % Modern package that handles PNG, JPG, PDF, EPS
\usepackage{subcaption}
\usepackage{calc}
\usepackage{amssymb}
\usepackage{amstext}
\usepackage{amsmath}
\usepackage{amsthm}
\usepackage{multicol}
\usepackage{pslatex}
\usepackage{apalike}
\usepackage{algorithm2e}  % Required by SCITEPRESS.sty even if not using algorithms
\usepackage[bottom]{footmisc}
\usepackage{url}
\usepackage{array}
\usepackage{multirow}
\usepackage{SCITEPRESS}  % Must be loaded AFTER algorithm2e

\begin{document}

\title{Self-Hosted High-Availability IAM System: A Reproducible Architecture}

\author{\authorname{Shahadat Hossain\sup{1}, Nuno G. Rodrigues\sup{1} and Rui P. Lopes\sup{2}}
\affiliation{\sup{1}Dept. of Informatics and Communications, ESTiG, Instituto Polit\'ecnico de Bragan\c{c}a, Bragan\c{c}a, Portugal}
\affiliation{\sup{2}Research Center in Digitalization and Intelligent Robotics (CeDRI), Instituto Polit\'ecnico de Bragan\c{c}a, Bragan\c{c}a, Portugal}
\email{hossain@ipb.pt, nuno@ipb.pt, rlopes@ipb.pt}
}

\keywords{Keycloak, Identity and Access Management, Single Sign-On, High Availability, HAProxy, Patroni, PostgreSQL, Docker Compose, Load Balancing, Database Failover, Infrastructure as Code, Open Source, Distributed Systems, Observability}

\abstract{Enterprise identity and access management (IAM) systems require high availability (HA) to prevent authentication failures that compromise security and disrupt operations. Organizations face critical deployment challenges: cloud-managed IAM services incur \$500--2000/month costs with elevated latencies, while vendor-specific HA solutions create unsuitable lock-in for regulatory compliance and multi-cloud strategies. Existing self-hosted approaches lack validated HA architectures with automated failover and empirical performance validation.
This paper presents a fully open-source, self-hosted HA architecture for Keycloak, integrating HAProxy load balancing, PostgreSQL with Patroni for automated failover, and comprehensive observability through Prometheus, Grafana, and Loki. Rigorous validation using the official Keycloak Gatling benchmark (n=10 trials) demonstrates sustained throughput of 180 req/s with sub-100ms latency and 0\% error rate, scaling to 362 req/s at saturation. Production-realistic failover testing under sustained load (n=3 trials, 181,500 requests) confirms 0\% error rate with sub-40ms mean latency and sub-30s recovery with zero data loss. Complete infrastructure-as-code artifacts---Ansible playbooks, Docker Compose files, and Grafana dashboards---enable reproducible one-command deployment on commodity hardware in under 10 minutes, achieving 61\% cost reduction versus AWS-managed equivalents. The validated architecture addresses a critical gap in production-grade Keycloak deployment guidance for mid-scale enterprises requiring infrastructure autonomy, regulatory compliance, or cost optimization.}

\onecolumn \maketitle \normalsize \vfill

\section{\uppercase{Introduction}}
\label{sec:introduction}

Identity and Access Management (IAM) systems serve as critical infrastructure for authentication and authorization in distributed enterprise applications (Grassi et al., 2017; Hardt, 2012). Production-grade IAM deployments must maintain high availability (HA) to prevent authentication failures that disrupt user access, compromise security posture, and violate service-level agreements (Samarati and De Capitani di Vimercati, 2001). Keycloak, an open-source IAM solution supporting OpenID Connect, OAuth 2.0, and SAML 2.0, has gained widespread adoption across industries (Keycloak Project, 2024). However, achieving production-grade HA remains challenging due to database clustering complexity and load balancer session affinity requirements.

Keycloak's official documentation recommends Amazon Aurora PostgreSQL for database HA (Keycloak Project, 2024b). While Aurora is a managed PostgreSQL-compatible service, it introduces AWS-specific dependencies, limiting portability for organisations requiring on-premises deployment, regulatory data residency compliance, or multi-cloud flexibility. This creates an unmet need for validated, reproducible self-hosted Keycloak HA architectures using lightweight orchestration tools.

Existing research addresses isolated aspects of Keycloak deployment. Performance evaluation studies (Bunn and Miers, 2024) analyze authentication protocol overhead but lack HA validation under realistic load. Security-focused research (Ruhur et al., 2025) examines threat models but does not address infrastructure resilience. Kubernetes-based approaches (Sousa et al., 2025) demonstrate elastic horizontal scaling (2100 req/s) but introduce overhead that may not be justified for fixed-capacity deployments where workload is predictable and operational simplicity is prioritised.

The novelty of this work lies in three dimensions: (1) engineering integration---combining HAProxy, Patroni, and Keepalived into a cohesive Docker Compose architecture without cloud dependencies; (2) empirical validation---using the official Keycloak Gatling benchmark with statistically rigorous multi-trial methodology; and (3) methodological reproducibility---publishing complete infrastructure-as-code artifacts enabling independent replication on commodity hardware.

This paper makes three contributions: (1) a reproducible HA architecture combining HAProxy SSL/SNI routing with sticky sessions, Patroni-orchestrated PostgreSQL streaming replication, and Keepalived VIP management---eliminating cloud vendor dependencies; (2) empirical performance validation via the official Keycloak Gatling framework across 10 independent trials, demonstrating linear scalability (50--200~req/s, sub-100~ms latency), graceful saturation (362.1±3.3~req/s, 0\% errors), and production-realistic failover under 200~req/s sustained load (n=3 trials, 0\% error rate, 38±9~ms mean latency, 25--30~s recovery); and (3) complete infrastructure-as-code artifacts via Ansible playbooks enabling one-command deployment, publicly available at (Hossain et al., 2025).

\section{\uppercase{Related Work}}
\label{sec:related}

PostgreSQL HA solutions employ streaming replication with automated failover (PostgreSQL Global Development Group, 2023). Patroni (Patroni Contributors, 2023) leverages distributed consensus (etcd Authors, 2023) for leader election, achieving sub-30~second recovery. Load balancing strategies for stateful applications require session affinity (Cardellini et al., 1999; HAProxy Technologies, 2023), with HAProxy providing robust SSL/SNI routing capabilities (Tarreau, 2023). Bunn and Miers (2024) analyzed OpenID Connect protocol overhead in Keycloak, identifying token generation as a primary latency contributor (40--60~ms baseline), but lacked HA configuration and concurrent user simulation. Ruhur et al. (2025) evaluated Keycloak security patterns in microservices but did not address infrastructure resilience.

Sousa et al. (2025) demonstrated Kubernetes-orchestrated Keycloak achieving 2100~req/s with horizontal pod auto-scaling. Container orchestration platforms provide elastic scaling but introduce operational complexity (Casalicchio and Perciballi, 2017) and recovery times exceeding 60~seconds. Our Docker Compose-based architecture achieves comparable failover performance (25--30~s) with superior per-request latency.

Managed IAM services (Okta, Auth0, Azure AD) provide turnkey HA but incur high recurring costs, with enterprise plans ranging from \$500--2000/month depending on user volume and feature requirements (Identity Management Institute, 2024; Amazon Web Services, 2024). Cloud-based IAM services often exhibit elevated latencies due to multi-tenant architecture and geographic distribution (Bermbach et al., 2013). Keycloak's AWS Aurora pattern (Keycloak Project, 2024b) reduces infrastructure costs but maintains vendor lock-in. Self-hosted approaches offer cost advantages (Li et al., 2010), but existing guides lack validated architectures with empirical performance data. Docker Swarm provides native clustering but lacks first-class stateful database failover orchestration; Patroni was selected for its purpose-built PostgreSQL HA with etcd-based consensus, a capability absent from generic container orchestrators.

\section{\uppercase{System Architecture}}
\label{sec:architecture}

The proposed HA architecture employs a three-tier design separating database persistence, application logic, and traffic routing concerns (Figure~\ref{fig:system-architecture}).

\begin{figure*}[t]
\centering
\includegraphics[width=0.95\textwidth]{System Design HA.png}
\caption{High-availability system architecture with segregated database VIP (172.29.65.100) for internal cluster communication and public VIP (193.136.194.100) for external client access. Dashed lines indicate VRRP and database failover paths; solid lines represent active data flows.}
\label{fig:system-architecture}
\end{figure*}

\subsection{Database Persistence Layer}

Three PostgreSQL 16.4 nodes operate in streaming replication managed by Patroni 3.3.2. PostgreSQL was selected as the exclusive database layer due to its mature streaming replication protocol, native support by Patroni for consensus-based leader election, and the availability of robust open-source tooling (pg\_exporter, pgBouncer). While Keycloak also supports MariaDB, MySQL, MSSQL, and Oracle, none offer equivalent ecosystem maturity for automated HA without additional proprietary components. Synchronous replication (\texttt{synchronous\_standby\_names='ANY 1 (*)'}) guarantees zero data loss during failover by requiring write confirmation from at least one standby. The etcd 3.5 distributed key-value store provides consensus for leader election (Ongaro and Ousterhout, 2014) with 2-second health checks, enabling rapid failure detection. A floating Virtual IP (VIP) address (172.29.65.100) managed by Keepalived follows the Patroni leader, ensuring connections always target the active primary without DNS delays.

\subsection{Application Logic Layer}

Three Keycloak 26.3.2 containers form a clustered deployment using JDBC\_PING for dynamic node discovery via shared database queries. Each instance connects to the database VIP (172.29.65.100), ensuring automatic failover transparency. The cluster configures \texttt{KC\_PROXY=edge} to trust X-Forwarded-* headers and \texttt{KC\_HOSTNAME\_STRICT=false} for VIP access. JDBC\_PING eliminates multicast needs by querying a shared database table (\texttt{JGROUPSPING}) for active cluster members.

\subsection{Traffic Routing Layer}

Two bastion nodes host HAProxy 2.8 instances providing SSL termination, SNI-based routing, and sticky session management. Port 443 inspects the TLS Client Hello Server Name Indication (SNI) field to route traffic without decryption overhead. Sticky sessions use \texttt{cookie SERVERID insert indirect nocache} with \texttt{leastconn} algorithm to prevent authentication redirect loops. Health checks issue HTTP GET to \texttt{/health/ready} every 2~seconds. Keepalived manages a public VIP (193.136.194.100) using VRRP (Virtual Router Redundancy Protocol) (Nadas, 2010) with priority-based active-passive failover, enabling 2--3~second detection.

\subsection{Observability Infrastructure}

Prometheus 2.54 (Prometheus Authors, 2023) scrapes metrics from HAProxy exporter, postgres\_exporter, Patroni, and Keycloak with 15-day retention. Grafana 11.3 (Grafana Labs, 2023a) visualizes metrics through pre-configured dashboards. Loki 3.2 (Grafana Labs, 2023b) aggregates logs from Promtail agents with Docker socket access, enabling log correlation by container labels for root-cause analysis (Turnbull, 2014). This observability triad follows modern monitoring best practices (Beyer et al., 2016), enabling sub-5~second alert detection during failures.

\subsection{Deployment Automation}

Ansible 2.17 orchestrates reproducible multi-node provisioning through the playbooks: backend cluster initialization, database VIP management, PostgreSQL access control, monitoring user provisioning, HAProxy deployment, metric exporter installation, and observability stack deployment. All services deploy via Docker Compose 2.29, enabling versioned, reproducible deployments and straightforward rollback by re-deploying a prior configuration. While Docker Compose does not provide true atomic rolling updates---which require additional orchestration tooling such as Docker Swarm or Kubernetes---it suffices for the controlled, sequential update procedures appropriate to fixed-capacity deployments. The automation achieves one-command deployment in 8--10~minutes on vanilla Debian 12 nodes, maximising reproducibility across commodity hardware without pre-configured base images.

\section{\uppercase{Experimental Methodology}}
\label{sec:methodology}

\subsection{Testbed Configuration}

The experimental validation employs a five-node testbed architecture running Debian 12, interconnected via a dedicated 1~Gbps private network. The testbed design separates database, application, and routing tiers to enable independent component analysis and realistic failure scenarios.

Three backend nodes (4 vCPU, 8GB RAM, 40GB SSD each, running Debian 12) form the core infrastructure. Each node runs an identical software stack: PostgreSQL 16.4 (data directory \texttt{/data/patroni}), Patroni 3.3.2 (configuration at \texttt{/etc/patroni/patroni.yml}), etcd 3.5 (peer port 2380, client port 2379), and Keycloak 26.3.2 (HTTP port 8080, bound to internal network interface). All three nodes participate equally in Patroni leader election; any node may be promoted to primary without configuration changes. Complete Ansible playbooks and Docker Compose files reproducing this configuration are available at (Hossain et al., 2025).

The routing tier comprises two bastion nodes, which provide load balancing and observability. Bastion1 (8 vCPU, 8GB RAM) serves as the primary entry point, hosting HAProxy, Keepalived, and the complete observability stack (Prometheus, Grafana, Loki); the higher vCPU count accommodates SSL termination overhead. Bastion2 (4 vCPU, 8GB RAM) operates in standby mode with HAProxy and Keepalived, enabling transparent failover of traffic routing functions.

\subsection{Load Testing Framework}

We employ the official Keycloak Gatling benchmark (Keycloak Project, 2023) with \texttt{ClientSecret} scenario simulating OAuth 2.0 client credentials flow, using Gatling's Open Workload Model (Gatling Corp, 2023) to maintain constant arrival rates and avoid coordinated omission (Thompson, 2015). Each test executes a 5-second ramp-up followed by a 30-second sustained load at target rates (50, 100, 200, 400~req/s) with 10 independent trials per level (Jain, 1991; Georges et al., 2007). Figure~\ref{fig:response-time-distribution} shows a representative trial; Table~\ref{tab:gatling-results} reports aggregated statistics.

\subsection{Failover Testing Protocol}

Database failover validation employs controlled failure injection. We verify cluster health via Patroni REST API, then execute \texttt{docker stop patroni-backend1} to simulate primary failure while monitoring logs for leader election. Time-to-recovery measures from failure detection to new leader promotion. VIP migration is validated via \texttt{ip addr show} and HAProxy health checks. Zero data loss is confirmed by comparing PostgreSQL transaction IDs pre- and post-failover.

\section{\uppercase{Results and Analysis}}
\label{sec:results}

\subsection{Authentication Performance Under Progressive Load}

Table~\ref{tab:gatling-results} presents empirical performance metrics across four load levels. The system exhibits linear scalability from 50--200~req/s with stable latency: mean latency increases only 4.7\% (85~ms to 89~ms) despite 4× throughput growth.

\begin{table}[!t]
\centering
\caption{Keycloak Authentication Performance Under OAuth 2.0 Client Credentials Load}
\label{tab:gatling-results}
\renewcommand{\arraystretch}{1.2}
\scriptsize
\begin{tabular}{|c|c|c|c|c|c|c|}
\hline
\textbf{Target} & \textbf{Actual} & \textbf{Mean} & \textbf{P50} & \textbf{P95} & \textbf{P99} & \textbf{Err} \\
\textbf{(r/s)} & \textbf{(r/s)} & \textbf{(ms)} & \textbf{(ms)} & \textbf{(ms)} & \textbf{(ms)} & \textbf{(\%)} \\
\hline
50 & 45 & 85 & 82 & 85 & 124 & 0 \\
\hline
100 & 90 & 85 & 83 & 86 & 109 & 0 \\
\hline
200 & 180 & 89 & 86 & 98 & 143 & 0 \\
\hline
400 & 362±3 & 67±55 & 33±2 & 206±353 & 720±974 & 0 \\
\hline
\multicolumn{7}{l}{\scriptsize Statistical validation: n=10 trials, 30s window, 8-core bastion} \\
\end{tabular}
\end{table}

At 400~req/s, 10-trial validation yields 362.1±3.3~req/s (CV=0.9\%), median 32.9±1.9~ms (95\% CI: 31.6--34.2~ms), and 0\% error rate across 130,000 requests.

Figure~\ref{fig:response-time-distribution} displays response time distribution from a representative trial, showing tight clustering where 50\% complete under 31~ms (median), 75\% under 35~ms, and 95\% under 51~ms.

\begin{figure}[!t]
\centering
\includegraphics[width=\columnwidth]{gatling-response-time-distribution.png}
\caption{Response time distribution under maximum stress test (Trial 1, n=10 validation set): 86\% of requests complete within 25--35~ms; median 31~ms, P95 51~ms, P99 89~ms. Aggregated across all trials: median 32.9±1.9~ms (95\% CI: 31.6--34.2~ms). Note: x-axis is linear scale.}
\label{fig:response-time-distribution}
\end{figure}

\subsection{Comparative Performance Analysis}

\textbf{Cloud IAM Services:} Mean latency of 67.3±54.9~ms and median 32.9±1.9~ms demonstrate single-tenant advantages over multi-tenant cloud IAM services exhibiting noisy-neighbor effects (Bermbach et al., 2013).

\textbf{Kubernetes-Orchestrated Deployments:} Sousa et al. (2025) achieved 2100~req/s with elastic scaling, approximately 6× our 362~req/s, reflecting different objectives: Kubernetes suits variable workloads, while our fixed cluster prioritises operational simplicity for stable loads. Our 0\% failover error rate contrasts with their 0.4\% baseline, attributable to HAProxy cookie-based sticky sessions.

\subsection{Database Failover Resilience}

Controlled PostgreSQL primary termination initiated Patroni's automated failover: etcd heartbeat miss triggered Raft consensus voting (Ongaro and Ousterhout, 2014) (t=0--5s), Keepalived reassigned the database VIP (t=5--15s), and Keycloak connection pools fully reconnected by t=20s. Total recovery measured 20--30 seconds. Grafana monitoring (Figure~\ref{fig:grafana-dashboard}) confirms zero transaction loss via contiguous PostgreSQL transaction IDs.

\begin{figure*}[!t]
\centering
\includegraphics[width=0.95\textwidth]{grafana-dashboard.png}
\caption{Grafana monitoring dashboard during database failover event showing (top) PostgreSQL replication lag spike during leader election, (middle) Patroni role transition showing backend1 termination and backend3 promotion to leader, and (bottom) PostgreSQL transaction continuity proving zero data loss through uninterrupted commit progression.}
\label{fig:grafana-dashboard}
\end{figure*}

During failover, Keycloak experienced transient queuing but zero hard failures. HAProxy's \texttt{retries 3} configuration automatically resubmits requests encountering database connection errors, transparent to end-users.

\subsection{Failover Resilience Under Sustained Load}

Three independent trials executed a 200~req/s load for 5 minutes with controlled Patroni leader termination at t=120~seconds. Statistical validation across 181,500 total requests demonstrated consistent failover resilience (CV\textless24\%).

Table~\ref{tab:failover-under-load} presents aggregated metrics. The system maintained 0\% error rate in all trials, confirming HAProxy's retry mechanism (\texttt{retries 3}) effectively masked database unavailability. Mean latency was 38±9~ms with median 27±6~ms; P99 of 191±46~ms remained within typical sub-second SLA requirements (Bertino and Takahashi, 2010).

\begin{table}[!t]
\centering
\caption{Database failover performance under sustained 200~req/s load}
\label{tab:failover-under-load}
\small
\begin{tabular}{|l|c|c|}
\hline
\textbf{Metric} & \textbf{Mean±SD} & \textbf{Range} \\
\hline
Mean Latency (ms) & 38±9 & 32--48 \\
Median (P50, ms) & 27±6 & 24--34 \\
P95 Latency (ms) & 83±15 & 70--99 \\
P99 Latency (ms) & 191±46 & 138--223 \\
Max Spike (ms) & 2061±9 & 2055--2071 \\
Error Rate (\%) & 0.0 & 0.0 \\
Throughput (req/s) & 197.9±0.4 & 197.7--198.4 \\
\hline
\multicolumn{3}{l}{\scriptsize Statistical validation: n=3 independent trials, 60,500 requests/trial} \\
\end{tabular}
\end{table}

Transient spikes to 2061±9~ms affected only 0.22\% of requests; 99.78\% completed under 800~ms. Low variance (CV\textless5\%) confirms system stability under failure. Our 0\% error rate during active failover validates that HAProxy retry logic and Patroni leader election enable transparent failover for stateless authentication workloads.

\subsection{HAProxy VIP Failover Performance}

Simulated bastion1 failure triggered Keepalived VRRP failover: bastion2 detected a missed advertisement after 3~seconds, claimed VIP (193.136.194.100), and broadcast gratuitous ARP. ICMP confirmed 2--3~second interruption; sticky sessions required no client reconnection. Across all 130,000+ requests, zero \texttt{ERR\_TOO\_MANY\_REDIRECTS} errors occurred, with even backend distribution: backend1 (33.2\%), backend2 (33.5\%), backend3 (33.3\%).

\subsection{Observability and Incident Detection}

Prometheus detected database failover via the \texttt{patroni\_postgres\_running} gauge transitioning from 1 to 0 on backend1, triggering an alert within 3~seconds. All 15 configured metric exporters reported UP status during steady-state operations.

Loki log aggregation correlated time-aligned events across Patroni, PostgreSQL, Keepalived, and HAProxy logs, isolating failover events within a 3-second window for rapid incident diagnosis. Log timestamps confirmed the recovery sequence: leader election at t+2s, PostgreSQL shutdown at t+3s, and VIP reassignment at t+5s, consistent with the measured 25--30~second total recovery time. This multi-signal observability reduces mean time to detection (MTTD) and mean time to recovery (MTTR) through comprehensive situational awareness.

\section{\uppercase{Discussion}}
\label{sec:discussion}

\subsection{Economic Analysis}

Table~\ref{tab:cost-analysis} presents a comparative cost analysis over 12 months.

\begin{table}[!t]
\centering
\caption{Annual Cost Comparison: Self-Hosted vs. AWS-Managed Deployment}
\label{tab:cost-analysis}
\small
\begin{tabular}{|l|r|r|}
\hline
\textbf{Component} & \textbf{AWS} & \textbf{Self-Hosted} \\
\hline
Database (Aurora Multi-AZ) & \$5,400 & --- \\
Application Load Balancer & \$1,200 & --- \\
EC2 Instances (3× t3.xlarge) & \$2,628 & --- \\
\hline
Private Cloud VMs (5× 4vCPU) & --- & \$2,400 \\
SSL Certificates (Let's Encrypt) & --- & \$0 \\
Operational Labor (20h/month) & --- & \$1,200 \\
\hline
\textbf{Total Annual Cost} & \textbf{\$9,228} & \textbf{\$3,600} \\
\textbf{Cost Savings} & --- & \textbf{61\%} \\
\hline
\end{tabular}
\end{table}

Self-hosted deployment achieves 61\% cost reduction (\$5,628 annual savings). \textbf{Cost model assumptions:} AWS pricing as of Q4 2024; single-region deployment; operational labour at \$40/hour for 20 hours/month. Intangible costs---Patroni tuning, Ansible maintenance, and 24/7 incident response---are not included and may reduce savings for organisations without existing DevOps capability.

\subsection{Operational Tuning Recommendations}

Production deployment revealed critical tuning parameters:

\textbf{Patroni:} Reduce \texttt{loop\_wait} from 10s to 2s to improve failover responsiveness by 60\% (40--50s~→~25--30s); set \texttt{retry\_timeout: 10}, \texttt{ttl: 30}.

\textbf{HAProxy:} Match \texttt{nbthread} to vCPU count; the 4-core configuration bottlenecked at 367\% CPU (493ms latency), while 8 cores reduced mean latency to 67ms.

\textbf{PostgreSQL Pooling:} Set \texttt{KC\_DB\_POOL\_MAX\_SIZE: 30} per node (90 total) to avoid exhausting \texttt{max\_connections=200}.

\subsection{Limitations and Future Research}

This evaluation addresses single-datacenter HA and does not examine multi-region replication. Key limitations include:

\textbf{Statistical Validation:} While 10-trial validation for performance characterization and 3-trial validation for failover testing provide robust baseline metrics (CV\textless5\%), extending to larger sample sizes (n=30) would strengthen confidence intervals, particularly for tail latencies exhibiting high variability.

\textbf{Security Hardening:} Security hardening is intentionally scoped out to focus on HA validation. The architecture supports TLS 1.3 via HAProxy SSL termination using certificates from any Certificate Authority. The infrastructure-as-code artifacts include automated Let's Encrypt certificate provisioning for public reproducibility, while supporting organisational CA certificates for enterprise environments; production deployment uses IPB-issued certificates on the ipb.pt domain. Full hardening---PostgreSQL mutual TLS, admin console isolation, secrets management via HashiCorp Vault, and automated patching---is deferred, with penetration testing and compliance certification as planned extensions.

\textbf{Distributed Tracing:} The current observability stack omits distributed tracing. Integrating OpenTelemetry with Jaeger or Tempo would enable end-to-end request tracing across HAProxy, Keycloak, and PostgreSQL for deeper latency attribution---a planned extension.

\textbf{Disaster Recovery:} Complete datacenter failure requires backup procedures. Future work should validate PostgreSQL WAL archival, point-in-time recovery (RPO \textless1 hour), and cross-datacenter asynchronous replication for multi-region deployments.

\section{\uppercase{Conclusion}}
\label{sec:conclusion}

This paper presented a reproducible, self-hosted high-availability architecture for Keycloak, integrating HAProxy, Patroni, and comprehensive observability. Empirical validation with statistical rigor (n=10 trials) demonstrated sustained 180~req/s throughput with 89~ms mean latency under 200 concurrent users, scaling to 362.1±3.3~req/s at saturation with maintained reliability. Production-realistic failover testing under sustained 200~req/s load (n=3 trials, 181,500 requests) confirmed 0\% error rate with 38±9~ms mean latency during 25--30~second recovery, validating transparent failover for stateless authentication workloads.

The architecture addresses critical gaps in production-grade Keycloak deployment guidance by providing cloud-agnostic infrastructure suitable for mid-scale enterprises (150--300 concurrent users) requiring predictable performance without Kubernetes complexity. Complete infrastructure-as-code artifacts are publicly available at (Hossain et al., 2025).

Future research directions include: (1) multi-region replication and geographic distribution; (2) comprehensive security hardening including TLS mutual authentication, secrets management, and compliance certification; (3) extended failure scenario testing covering network partitions and simultaneous multi-node failures; (4) extended statistical validation (n=30) to tighten tail-latency confidence intervals; (5) mixed workload evaluation (authorisation code flow, token refresh, user federation); (6) cloud provider comparison extended to GCP and Azure, and direct empirical comparison with Kubernetes-based Keycloak deployments under identical hardware and workload conditions to quantify performance and operational complexity trade-offs; and (7) distributed tracing via OpenTelemetry and Jaeger. Organizations seeking infrastructure autonomy, regulatory compliance, or cost optimization can leverage these artifacts to deploy production-grade IAM without managed service dependencies.

\section*{ACKNOWLEDGEMENTS}

\begin{sloppypar}
This work was supported by the Instituto Polit\'ecnico de Bragan\c{c}a and FCT--\allowbreak{}Funda\c{c}\~ao para a Ci\^encia e a Tecnologia, I.P., by projects: CeDRI, UID/05757/2025 (DOI: \url{10.54499/UID/05757/2025}) and UID/PRR/\allowbreak{}05757/2025 (DOI: \url{10.54499/UID/PRR/05757/2025}); SusTEC, LA/P/0007/2020 (DOI: \url{10.54499/LA/P/0007/2020}).
\end{sloppypar}

The authors used Claude (Anthropic) to assist with language clarity and grammatical accuracy in portions of this manuscript. All content was reviewed, verified, and edited by the authors, who take full responsibility for the manuscript's accuracy and integrity.

\bibliographystyle{apalike}
\begin{thebibliography}{}

\bibitem[Amazon Web Services, 2024]{awspricing2024}
Amazon Web Services (2024).
\newblock Amazon RDS for PostgreSQL Pricing.
\newblock AWS Official Documentation.
\newblock Available at: \url{https://aws.amazon.com/rds/postgresql/pricing/}.

\bibitem[Bermbach et al., 2013]{bermbach2013cloud}
Bermbach, D., Kuhlenkamp, J., Derre, B., Klems, M., and Tai, S. (2013).
\newblock A Middleware Guaranteeing Client-Centric Consistency on Top of Eventually Consistent Datastores.
\newblock In {\em Proc. IEEE Int. Conf. Cloud Engineering (IC2E)}, pages 114--123.

\bibitem[Bertino and Takahashi, 2010]{bertino2010identity}
Bertino, E. and Takahashi, K. (2010).
\newblock {\em Identity Management: Concepts, Technologies, and Systems}.
\newblock Artech House.

\bibitem[Beyer et al., 2016]{beyer2016sre}
Beyer, B., Jones, C., Petoff, J., and Murphy, N.~R. (2016).
\newblock {\em Site Reliability Engineering: How Google Runs Production Systems}.
\newblock O'Reilly Media.

\bibitem[Bunn and Miers, 2024]{bunn2024}
Bunn, C. D.~S. and Miers, C.~C. (2024).
\newblock Evaluating Performance Impacts in Identity Management Based on Keycloak and OpenID Connect.
\newblock In {\em Proc. Brazilian Symp. Information and Computational Systems Security (SBSeg)}.

\bibitem[Burns et al., 2016]{burns2016kubernetes}
Burns, B., Grant, B., Oppenheimer, D., Brewer, E., and Wilkes, J. (2016).
\newblock Borg, Omega, and Kubernetes.
\newblock {\em ACM Queue}, 14(1):70--93.

\bibitem[Cardellini et al., 1999]{cardellini2002dynamic}
Cardellini, V., Colajanni, M., and Yu, P.~S. (1999).
\newblock Dynamic Load Balancing on Web-Server Systems.
\newblock {\em IEEE Internet Computing}, 3(3):28--39.

\bibitem[Casalicchio and Perciballi, 2017]{casalicchio2019container}
Casalicchio, E. and Perciballi, V. (2017).
\newblock Auto-Scaling of Containers: The Impact of Relative and Absolute Metrics.
\newblock In {\em Proc. IEEE Int. Conf. Fog and Edge Computing}, pages 207--214.

\bibitem[Docker Inc., 2023]{dockercompose2023}
Docker Inc. (2023).
\newblock Compose Specification.
\newblock Docker Documentation.
\newblock Available at: \url{https://docs.docker.com/compose/}.

\bibitem[etcd Authors, 2023]{etcd2023}
etcd Authors (2023).
\newblock etcd: A Distributed, Reliable Key-Value Store for the Most Critical Data of a Distributed System.
\newblock Available at: \url{https://etcd.io}.

\bibitem[Gatling Corp, 2023]{gatling2023}
Gatling Corp (2023).
\newblock Open vs. Closed Workload Models: Understanding Injection Patterns.
\newblock Gatling Documentation.
\newblock Available at: \url{https://gatling.io/docs/gatling/reference/current/core/injection/}.

\bibitem[Georges et al., 2007]{georges2007statistically}
Georges, A., Buytaert, D., and Eeckhout, L. (2007).
\newblock Statistically Rigorous Java Performance Evaluation.
\newblock In {\em Proc. ACM SIGPLAN Conf. Object-Oriented Programming Systems and Applications (OOPSLA)}, pages 57--76.

\bibitem[Grafana Labs, 2023a]{grafana2023}
Grafana Labs (2023a).
\newblock Grafana: The Open Observability Platform.
\newblock Available at: \url{https://grafana.com/}.

\bibitem[Grafana Labs, 2023b]{loki2023}
Grafana Labs (2023b).
\newblock Loki: Like Prometheus, but for Logs.
\newblock Available at: \url{https://grafana.com/oss/loki/}.

\bibitem[Grassi et al., 2017]{grassi2017}
Grassi, P.~A., Garcia, M.~E., and Fenton, J.~L. (2017).
\newblock Digital Identity Guidelines: Authentication and Lifecycle Management.
\newblock NIST Special Publication 800-63B.

\bibitem[HAProxy Technologies, 2023]{haproxydocs2023}
HAProxy Technologies (2023).
\newblock HAProxy Configuration Manual.
\newblock Available at: \url{https://www.haproxy.org/}.

\bibitem[Hardt, 2012]{hardt2012oauth}
Hardt, D. (2012).
\newblock The OAuth 2.0 Authorization Framework.
\newblock RFC 6749, Internet Engineering Task Force.

\bibitem[Identity Management Institute, 2024]{openiamcost2024}
Identity Management Institute (2024).
\newblock Total Cost of Ownership for Enterprise Identity Management: On-Premises vs Cloud.
\newblock White Paper.
\newblock Available at: \url{https://www.identitymanagementinstitute.org/}.

\bibitem[Jain, 1991]{jain1991art}
Jain, R. (1991).
\newblock {\em The Art of Computer Systems Performance Analysis}.
\newblock John Wiley \& Sons.

\bibitem[Keycloak Project, 2023]{keycloakBenchmark2023}
Keycloak Project (2023).
\newblock keycloak-benchmark: Official Load Testing and Performance Benchmarking Framework.
\newblock GitHub.
\newblock Available at: \url{https://github.com/keycloak/keycloak-benchmark}.

\bibitem[Keycloak Project, 2024a]{keycloakdocs2024}
Keycloak Project (2024a).
\newblock Keycloak Documentation.
\newblock Available at: \url{https://www.keycloak.org/documentation}.

\bibitem[Keycloak Project, 2024b]{keycloakAurora2024}
Keycloak Project (2024b).
\newblock Deploying AWS Aurora in Multiple Availability Zones for High Availability.
\newblock Keycloak Documentation.
\newblock Available at: \url{https://www.keycloak.org/high-availability/deploy-aurora-multi-az}.

\bibitem[Li et al., 2010]{li2010cloudcmp}
Li, A., Yang, X., Kandula, S., and Zhang, M. (2010).
\newblock CloudCmp: Comparing Public Cloud Providers.
\newblock In {\em Proc. ACM SIGCOMM Internet Measurement Conf. (IMC)}, pages 1--14.

\bibitem[Merkel, 2014]{merkel2014docker}
Merkel, D. (2014).
\newblock Docker: Lightweight Linux Containers for Consistent Development and Deployment.
\newblock {\em Linux Journal}, 2014(239).

\bibitem[Nadas, 2010]{rfc5798vrrp}
Nadas, S. (2010).
\newblock Virtual Router Redundancy Protocol (VRRP) Version 3 for IPv4 and IPv6.
\newblock RFC 5798, Internet Engineering Task Force.

\bibitem[Ongaro and Ousterhout, 2014]{ongaro2014raft}
Ongaro, D. and Ousterhout, J. (2014).
\newblock In Search of an Understandable Consensus Algorithm.
\newblock {\em USENIX Annual Technical Conference}, pages 305--320.

\bibitem[Patroni Contributors, 2023]{patroni2023}
Patroni Contributors (2023).
\newblock Patroni: A Template for PostgreSQL High Availability with ZooKeeper, etcd, Consul, or Kubernetes.
\newblock GitHub.
\newblock Available at: \url{https://github.com/patroni/patroni}.

\bibitem[PostgreSQL Global Development Group, 2023]{postgres2023replication}
PostgreSQL Global Development Group (2023).
\newblock High Availability, Load Balancing, and Replication.
\newblock PostgreSQL Documentation.
\newblock Available at: \url{https://www.postgresql.org/docs/current/high-availability.html}.

\bibitem[Prometheus Authors, 2023]{prometheus2023}
Prometheus Authors (2023).
\newblock Prometheus: Monitoring System and Time Series Database.
\newblock Available at: \url{https://prometheus.io/}.

\bibitem[Ruhur et al., 2025]{ruhur2025}
Ruhur, W., Girsang, A.~S., and Warnars, H. L. H.~S. (2025).
\newblock Design and Security Evaluation of IAM Module in Microservice Architecture Using Keycloak.
\newblock {\em Indonesian Journal of Computing and Cybernetics Systems}, 14(2):145--156.

\bibitem[Samarati and De Capitani di Vimercati, 2001]{samarati2001access}
Samarati, P. and De Capitani di Vimercati, S. (2001).
\newblock Access Control: Policies, Models, and Mechanisms.
\newblock In {\em Foundations of Security Analysis and Design}, pages 137--196. Springer.

\bibitem[Sousa et al., 2025]{sousa2025}
Sousa, A., Branco, F., Reis, A., and Reis, M. J. C.~S. (2025).
\newblock A Container-Native IAM Framework for Secure Green Mobility Services.
\newblock {\em Information}, 16(9):802.

\bibitem[Tarreau, 2023]{tarreau2023haproxy}
Tarreau, W. (2023).
\newblock HAProxy: The Reliable, High Performance TCP/HTTP Load Balancer.
\newblock HAProxy Technologies.

\bibitem[Thompson, 2015]{thompson2015coordinated}
Thompson, G. (2015).
\newblock Coordinated Omission: A Major Source of Error in Latency Measurement.
\newblock Available at: \url{https://www.scylladb.com/2021/04/22/on-coordinated-omission/}.

\bibitem[Turnbull, 2014]{turnbull2014art}
Turnbull, J. (2014).
\newblock {\em The Art of Monitoring}.
\newblock CreateSpace Independent Publishing.

\bibitem[Hossain et al., 2025]{hossain2025github}
Hossain, S., Rodrigues, N.~G., and Lopes, R.~P. (2025).
\newblock Containerized High-Availability: Self-Hosted Keycloak HA with HAProxy, Patroni, and Ansible.
\newblock GitHub Repository.
\newblock Available at: \url{https://github.com/Shahadat-ngup/Containerized-HighAvailability}.

\bibitem[H\"{a}rder and Reuter, 1983]{haerder1983}
H\"{a}rder, T. and Reuter, A. (1983).
\newblock Principles of Transaction-Oriented Database Recovery.
\newblock {\em ACM Computing Surveys}, 15(4):287--317.

\bibitem[Bailis and Kingsbury, 2014]{bailis2014network}
Bailis, P. and Kingsbury, K. (2014).
\newblock The Network is Reliable.
\newblock {\em ACM Queue}, 12(7):20--32.

\bibitem[Kleppmann, 2017]{kleppmann2017designing}
Kleppmann, M. (2017).
\newblock {\em Designing Data-Intensive Applications}.
\newblock O'Reilly Media.

\end{thebibliography}

\end{document}